\documentclass[12pt]{article}
 
\usepackage[margin=1in]{geometry}
\usepackage{setspace}
\usepackage{times}
\usepackage{amsmath, amssymb}
\usepackage[hidelinks]{hyperref}
\usepackage{parskip}
 
\singlespacing

\begin{document}

\begin{titlepage}
    \centering
    \vspace*{2in}
    {\large MATH 447 -- Cryptography} \\[0.5cm]
    {\LARGE \textbf{Project Proposal}} \\[1.5cm]
    {\Large \textbf{End to End Encrypted Chat using ChaCha20-Poly1305}} \\[2cm]

    {\large Aidan, Maddie, Lei} \\[0.3cm]
    {\large April 2026}

    \vfill
\end{titlepage}

\section{Introduction}

    This proposal outlines the design and implementation of an end-to-end encrypted chat 
    application built in Python. The centrel cryptographic focus of this project is the 
    ChaCha20-Poly1305 encryption algorithm, which will be used to encrypt all messages 
    exchanged between clients. While the system also uses RSA-2048 for the initial 
    key exchange handshake, the primary subject of our study, analysis, and demonstration
    is the ChaCha20-Poly1305 cipher and its role in providing both confidentiality and  
    integrity for real time communication.

    ChaCha20-Poly1305 combines two primitives indepdently designed by Daniel J.\ 
    Bernstein, the ChaCha20 stream cipher (2008), which is a refinement of his 
    earlier Salsa20 cipher, and the Poly1305 one time message authentication code (2005).
    The combinbed authenticated encryption with associated data (AEAD) construction 
    was proposed for inclusion in IETF protocols in 2013 and was standardized in RFC~7539 (2015),
    later updated by RFC~8439 (2018) \cite{rfc8439}. Today, ChaCha20-Poly1305 is a 
    mandatory cipher suite in TLS~1.3 and is widely used in protocols such as OpenSSH,
    WireGuard, and IPsec, as well as in secure messaging applications like Signal \cite{rfc7905}.

    The motivation for studying ChaCha20-Poly1305 stems from its growing importance as an 
    alternative to AES-GCM. On devices that lack hardware AES acceleration, such as many ARM-based
    mobile processors, ChaCha20-Poly1305 typically delivers significantly better performance.
    Furthermore, ChaCha20 is designed to operate in constant time, making it resistant to timing based
    side channel attacks that can affect table lookup based implementations of AES \cite{cloudflare}.
    By building a working chat application that relies on this cipher, our project provides both a 
    theoretical exploration of the algorithm and a practical demonstration of its application in 
    a real world scenario.

\section{Mathematical and Cryptographic Background}

    \subsection{The ChaCha20 Stream Cipher}

        ChaCha20 is a stream cipher that generates a pseudorandom keystream by applying a series of add rotate XOR (ARX) operation to an internal state \cite{bernstein2008}. The cipher initializes a $4 \times 4$ matrix of 32 bit words (512 bits total) composed of four constant words, eight key words (256 bit key), one block counter word, and three nonce words (96 bit nonce). The core operation of ChaCha20 is the \textit{quarter round} function, which takes four 32 bit integers and applies a sequence of additions, XORs, and bitwise rotations. A complete round consists of four column quarter rounds followed by four diagonal quarter rounds. ChaCha20 applies 20 such rounds, hence the name, to the initial state, then adds the result to the original state to produce 64 bytes of keystream. The plaintext is encrypted by XORing it with this keystream, one 64 byte block at a time. The block counter is incremented for each subsequent block, allowing encryption of arbitraty length messages.

        The ARX design is signifigant from a cryptographic standpoint because it avoids the use of look up tables, which are common in AES. This means that ChaCha20 implementations can be made constant time without special effort, eliminating an entire class of timing vulnerabiliies. The cipher provides 256 bits of security, and no practical attacks against the full 20 round variant are known.

    \subsection{The Poly1305 Message Authentication Code}

        Poly1305 is a universal hash function that computes a 128 bit authentication tag over a message \cite{bernstein2005}. Given a 256 bit one time key consisting of two 128 bit halves, $r$ and $s$, the Poly1305 tag for a message is computed as a polynomial evaluation modulu the prime $p = 2^{130} - 5$. The message is divided into 16 byte blocks, where each block is treated as an integer with a high bit appended, and the running polynomial accumulator is updated by adding each block and multiplying by $r$. After all blocks are processed, the one time pad $s$ is added to the accumulator, and the result is taken module $2^{128}$ to produce the final 128 bit tag. The security of Poly1305 relies on the key being used only once. In the ChaCha20-Poly1305 construction, this requirement is satisfied automatically since the first 32 vytes of the ChaCha20 keystream are used as the one time Poly1305 key, with the remaining keystream used to encrypt the plaintext. This tight coupling between the cipher and the authenticator ensures that the Poly1305 key is always fresh and unique for each message, as long as the nonce is never reused with the same key. 

    \subsection{The AEAD Construction}

        Authenticated Encryption with Associated Data (AEAD) is a construction that provides both confidentiality and integrity in a single operation \cite{rfc8439}. The ChaCha20-Poly1305 AEAD takes as input a 256 bit key, a 96 bit nonce, a plaintext message, and optional associated data (AAD) that is authenticated but not encrypted. The encryption process generates a Poly1305 key from the first ChaCha20 block, encrypts the plaintext using the remainder of the keystream, and then computes an authentication tag over the AAD, the ciphertext, and their respective lengths. The resulting outpit is the ciphertext concatenated with the 128 bit tag. 

        Decryption reverses this process: the tag is verified before the plaintext is returned, ensuring that any tampering with the ciphertext or AAD is detected and rejected. This encrypt then MAC construction provides strong authenticity that avoids the well documented pit falls of cipher MAC combinations \cite{procter2014}. 

\section{Project Overview and Design}

    Our application follows a client server architecture with three components: A central relay server and two or more chat clients. The server facilitates message routing but is excluded from having access to the plaintext content of messages. All cryptographic operations occur on the client side, ensuring true end to end encryption. The project is implemented in Python using the \texttt{socket} library for networking and the \texttt{cryptography} library for all cryptographic primitives. 

    The communication protocol happens in two phases. In the first phase, clients exchange RSA-2048 public keys to perform an asymmetric key exchange. Each client generates an RSA key pair, shares its public key through the server, and uses the partner's public key to encrypt a shared symmetric key. This initial handshake established a 256 bit symmetric key known only to the communicating clients. In the second phase, all subsequent messages encrypted using ChaCha20-Poly1305 with a unique 96 bit nonce (generated randomly for each message), and the resulting ciphertext along with its authentication tag is sent through the server to the recipient. The server sees only the byte streams and cannot decrypt or forge messages.

    It is important to emphasize that while RSA-2048 is used as a practical way for establishing the shared key, it is not the focus of this project. Our study centers on the ChaCha20-Poly1305 algorithm, its mathematical foundations, its security properties, its performance characteristics, and its behavior in a live communication system.

\section{Project Goals}

    The project has three primary goals. First we aim to build a fully functional chat application where all messages are encrypted using ChaCha20-Poly1305, demonstrating the practical application of the cipher in a real system. Second, we will provide a detailed mathematical and theoretical analysis of the ChaCha20-Poly1305 construction, including the quarter round function, the polynomial arithmetic underlying Poly1305, and the security guarantees of the AEAD composition. Third, we will construct a side by side comparison demonstrating the difference between encrypted and unencrypted communication, allowing us to visually show that the server cannot read messages when ChaCha20-Poly1305 is employed.

    The final project will include the complete Python source code (server and client), a ten page paper detailing the theory, implementation, and analysis of ChaCha20-Poly1305, a live demonstration of the encrypted chat system, and an in class presentation covering the cryptographic concepts and our findings.

\section{Timeline}

    \textbf{Weeks 1--2 (Completed):} Implemented the basic server and client using Python sockets. Verified multi client connectivity and basic relay functionality.

    \textbf{Week 3--4:} Implement RSA-2048 and key exchange between clients. Verify that both clients derive the same shared symmetric key. Integrate ChaCha20-Poly1305 message encryption and decryption using the shared key. Test with multiple message exchanges and verify tag authentication.

    \textbf{Week 5:} Build the unencrypted version for comparison. Develop the demonstration showing server side visiblity of plaintext vs. ciphertext. Write the final paper and prepare the final presenation.

\begin{thebibliography}{9}

    \bibitem{bernstein2008}
        D.~J. Bernstein,
        ``ChaCha, a variant of Salsa20,''
        2008.
        Available: \url{https://cr.yp.to/chacha.html}

    \bibitem{bernstein2005}
        D.~J. Bernstein,
        ``The Poly1305-AES message-authentication code,''
        in \textit{Fast Software Encryption (FSE 2005)},
        Lecture Notes in Computer Science, vol.~3557,
        Springer, 2005, pp.~32--49.

    \bibitem{rfc8439}
        Y.~Nir and A.~Langley,
        ``ChaCha20 and Poly1305 for IETF Protocols,''
        RFC~8439, Internet Engineering Task Force, June 2018.
        Available: \url{https://www.rfc-editor.org/rfc/rfc8439}

    \bibitem{rfc7905}
        A.~Langley, W.~Chang, N.~Mayrogiannopoulos,
        J.~Strombergson, and S.~Josefsson,
        ``ChaCha20-Poly1305 Cipher Suites for Transport Layer
        Securtity (TLS),''
        RFC~7905, Internet Engineering Task Force, June 2016.
        Available: \url{https://www.rfc-editor.org/rfc/rfc7905}

    \bibitem{procter2014}
        G.~Procter,
        ``A Security Analysis of the Compisition of ChaCha20
        and Poly1305,''
        Cryptology ePrint Archive, Report 2014/613, 2014.
        Available: \url{https://eprint.iacr.org/2014/613}

\end{thebibliography}

\end{document}